\documentclass[11pt,a4paper]{article}

\usepackage[margin=1in]{geometry}
\usepackage{amsmath,amssymb,amsfonts,bm}
\usepackage{booktabs}
\usepackage{hyperref}
\usepackage{microtype}
\usepackage{enumitem}

\hypersetup{
  colorlinks=true,
  linkcolor=blue!60!black,
  urlcolor=blue!60!black,
  citecolor=blue!60!black
}

\title{Model 1: Two Coupled Classical Lorentz Oscillators\\
\large Analytical benchmark for PlasMol parallel absorption}
\author{PlasMol validation notes}
\date{\today}

\begin{document}
\maketitle

\begin{abstract}
\noindent
We treat the Au nanoparticle LSPR and the Na atomic transition as two
driven, damped harmonic oscillators coupled by a near-field constant
$\kappa$.  The linear system is solved in closed form for the
frequency-dependent displacements, yielding an effective molecular
polarizability $\alpha_m^{\mathrm{eff}}(\omega)=\mu_m(\omega)/E_{\mathrm{inc}}(\omega)$
that is directly comparable (up to peak normalization) to PlasMol's
hybrid Fourier spectrum
$A(\omega)\propto-\omega\,\mathrm{Im}[\mu(\omega)/E_{\mathrm{inc}}(\omega)]$.
\end{abstract}

\tableofcontents
\newpage

% =============================================================================
\section{Motivation and correspondence to PlasMol}
% =============================================================================

PlasMol's parallel Fourier path for a nanoparticle--molecule system
records the molecular induced dipole $\mu(t)$ under a broadband Meep
drive and divides by a vacuum incident field $E_{\mathrm{inc}}$ sampled
at the molecule site (no NP, no molecule).  The published spectrum is
\begin{equation}
  A(\omega)
  \;\propto\;
  -\omega\,
  \operatorname{Im}\!
  \Big[
    \frac{\mu(\omega)}{E_{\mathrm{inc}}(\omega)}
  \Big],
  \label{eq:plasmol-A}
\end{equation}
peak-normalized over a user energy window.

Equation~\eqref{eq:plasmol-A} is \emph{not} the hybrid extinction of the
composite.  It is the molecular channel relative to free-space
incidence, with nanoparticle scattering retained inside $\mu$.  Model~1
targets exactly that observable with the simplest two-level (two
oscillator) dynamics.

\paragraph{Target system (defaults in \texttt{spectrum\_model1.py}).}
\begin{center}
\begin{tabular}{@{}ll@{}}
\toprule
Subsystem & Parameters \\
\midrule
Au LSPR & $\omega_p = 2.384\,\mathrm{eV}$, damping $\gamma_p$ \\
Na transition & $\omega_m = 2.109\,\mathrm{eV}$, damping $\gamma_m$ \\
Detuning & $\delta = \omega_p - \omega_m \approx 0.275\,\mathrm{eV}$ \\
Coupling & $\kappa$ (units of energy$^2$), user-tunable \\
\bottomrule
\end{tabular}
\end{center}

% =============================================================================
\section{Bare Lorentz oscillators}
% =============================================================================

A single classical Lorentz oscillator with resonance $\omega_j$, damping
$\gamma_j$, and drive strength $f_j$ obeys
\begin{equation}
  \ddot x_j + \gamma_j \dot x_j + \omega_j^2 x_j
  =
  f_j\, E(t),
  \qquad
  j\in\{p,m\}.
  \label{eq:bare-eom}
\end{equation}
In the frequency domain ($e^{-i\omega t}$ convention, or equivalently
replacing $\partial_t\to -i\omega$ in the Fourier domain used here),
\begin{equation}
  \Omega_j^2(\omega)\, x_j(\omega)
  =
  f_j\, E(\omega),
  \qquad
  \Omega_j^2(\omega)
  \equiv
  \omega_j^2 - \omega^2 - i\gamma_j\omega.
  \label{eq:Omega}
\end{equation}
The bare polarizability (reduced units where $\mu_j \propto x_j$) is
\begin{equation}
  \alpha_j^{(0)}(\omega)
  =
  \frac{f_j}{\Omega_j^2(\omega)}
  =
  \frac{f_j}{\omega_j^2 - \omega^2 - i\gamma_j\omega}.
  \label{eq:alpha0}
\end{equation}
Absorption of an isolated oscillator is proportional to
$-\omega\,\operatorname{Im}\alpha_j^{(0)}(\omega)$, which is a
non-negative Lorentzian-like peak of width set by $\gamma_j$.

% =============================================================================
\section{Coupled equations of motion}
% =============================================================================

Near-field coupling between the plasmon and molecular oscillators is
encoded as a bilinear term with real coupling constant $\kappa$
(dimensions of frequency$^2$ / energy$^2$ when $\omega$ is in eV):
\begin{align}
  \Omega_p^2(\omega)\, x_p + \kappa\, x_m
  &=
  f_p\, E_{\mathrm{inc}},
  \label{eq:coup-p} \\
  \kappa\, x_p + \Omega_m^2(\omega)\, x_m
  &=
  f_m\, E_{\mathrm{inc}}.
  \label{eq:coup-m}
\end{align}
In matrix form,
\begin{equation}
  \begin{pmatrix}
    \Omega_p^2 & \kappa \\
    \kappa & \Omega_m^2
  \end{pmatrix}
  \begin{pmatrix}
    x_p \\ x_m
  \end{pmatrix}
  =
  \begin{pmatrix}
    f_p \\ f_m
  \end{pmatrix}
  E_{\mathrm{inc}}.
  \label{eq:matrix}
\end{equation}
The determinant
\begin{equation}
  \Delta(\omega)
  =
  \Omega_p^2\Omega_m^2 - \kappa^2
  \label{eq:det}
\end{equation}
vanishes at the hybrid poles.  Cramer's rule gives the closed-form
displacements
\begin{align}
  x_p(\omega)
  &=
  \frac{
    \Omega_m^2 f_p - \kappa f_m
  }{
    \Delta(\omega)
  }\,
  E_{\mathrm{inc}},
  \label{eq:xp} \\[0.4em]
  x_m(\omega)
  &=
  \frac{
    \Omega_p^2 f_m - \kappa f_p
  }{
    \Delta(\omega)
  }\,
  E_{\mathrm{inc}}.
  \label{eq:xm}
\end{align}

\paragraph{Connection to a coupling energy $g$.}
In the literature one often writes a Hamiltonian coupling
$\hbar g(a^\dagger\sigma_- + \mathrm{h.c.})$ or a force constant
proportional to $g$.  Near resonance, the level repulsion scale is set
by an effective $g$ with $g^2 \sim \kappa$ (order of magnitude, depending
on oscillator normalizations).  For detuning
$\delta=\omega_p-\omega_m$, clear mode hybridization (avoided crossing /
Rabi-like doublet in the linear spectrum) typically requires
$g \gtrsim |\delta|/2$ and $g$ comparable to the mean linewidth.
Weak $\kappa$ yields a single molecular peak that is shifted and
reshaped---the usual weak-coupling / Purcell-type regime.

% =============================================================================
\section{Effective molecular polarizability and absorption}
% =============================================================================

Matching PlasMol's definition, the molecular response relative to the
incident field is
\begin{equation}
  \alpha_m^{\mathrm{eff}}(\omega)
  \equiv
  \frac{\mu_m(\omega)}{E_{\mathrm{inc}}(\omega)}
  \;\propto\;
  \frac{x_m(\omega)}{E_{\mathrm{inc}}}
  =
  \frac{
    \Omega_p^2 f_m - \kappa f_p
  }{
    \Omega_p^2\Omega_m^2 - \kappa^2
  }.
  \label{eq:alpha-eff}
\end{equation}
The model spectrum is
\begin{equation}
  A(\omega)
  =
  -\omega\,
  \operatorname{Im}\!
  \big[
    \alpha_m^{\mathrm{eff}}(\omega)
  \big],
  \label{eq:A}
\end{equation}
subsequently peak-normalized,
\begin{equation}
  A_{\mathrm{norm}}(\omega)
  =
  \frac{A(\omega)}{\max_{\omega}|A(\omega)|},
\end{equation}
with a global sign flip if the dominant extremum is negative (same
practical convention as PlasMol's \texttt{orient\_spectrum\_sign}).

\paragraph{Limits.}
\begin{itemize}[leftmargin=1.4em]
  \item $\kappa\to 0$:
    $\alpha_m^{\mathrm{eff}}\to f_m/\Omega_m^2=\alpha_m^{(0)}$
    (bare Na).
  \item $f_m\to 0$ but $\kappa\neq 0$:
    the molecule can still be driven indirectly through the plasmon
    term $-\kappa f_p/\Delta$, i.e.\ a pure near-field imprint of the
    LSPR on $\mu_m$.
  \item Large $\kappa$: hybrid poles split; $A(\omega)$ can show two
    peaks even though only $\mu_m$ is plotted.
\end{itemize}

% =============================================================================
\section{Hybrid mode frequencies}
% =============================================================================

The undriven poles satisfy $\Delta(\omega)=0$.  A useful algebraic
estimate replaces $\Omega_j^2$ by complex constants
$z_j=\omega_j^2-i\gamma_j\omega_j$ evaluated at the bare peaks and solves
\begin{equation}
  (z_p-\lambda)(z_m-\lambda)-\kappa^2=0
  \quad\Rightarrow\quad
  \lambda_\pm
  =
  \frac{z_p+z_m}{2}
  \pm
  \frac12
  \sqrt{(z_p-z_m)^2+4\kappa^2},
\end{equation}
then $\omega_\pm\approx\sqrt{\lambda_\pm}$ (principal branch).  The
script prints $\omega_\pm$ as vertical guides.  They are indicative,
not exact roots of the full transcendental
$\Omega_p^2(\omega)\Omega_m^2(\omega)=\kappa^2$.

In the equal-damping, near-resonant textbook limit one recovers the
familiar form
\begin{equation}
  \omega_\pm
  \approx
  \frac{\omega_p+\omega_m}{2}
  -\frac{i}{4}(\gamma_p+\gamma_m)
  \pm
  \frac12
  \sqrt{
    \Big(
      \delta - \tfrac{i}{2}(\gamma_p-\gamma_m)
    \Big)^2
    + 4g^2
  },
\end{equation}
with $g$ the coupling energy related to $\kappa$ through the chosen
oscillator normalization.

% =============================================================================
\section{What to expect for Au (2.384 eV) + Na (2.109 eV)}
% =============================================================================

\begin{enumerate}[leftmargin=1.4em]
  \item \textbf{Weak $\kappa$.}
        Single peak near $\omega_m$, possibly red- or blue-shifted by a
        fraction of $\gamma_m$.  A weak shoulder near $\omega_p$ may
        appear if $f_p$ and $\kappa$ allow plasmon-mediated driving of
        $x_m$.
  \item \textbf{Moderate $\kappa$.}
        Noticeable shift of the molecular maximum; lineshape becomes
        asymmetric (Fano-like) when the broad plasmon continuum
        interferes with the narrow molecular resonance.
  \item \textbf{Strong $\kappa$.}
        Two hybrid peaks; energy splitting grows with $\sqrt{\kappa}$.
        Peak-normalization still forces $\max|A|=1$, so compare
        \emph{positions and relative heights}, not absolute scale.
\end{enumerate}

\noindent
Your PlasMol \texttt{parallel\_abs} result---one dominant peak near the
Na energy---is the weak-to-moderate coupling pattern.  Model~1 validates
the \emph{method} when:
\begin{itemize}[leftmargin=1.4em]
  \item $\kappa=0$ recovers the bare Na lineshape used as input;
  \item increasing $\kappa$ shifts / splits $A(\omega)$ in the direction
        predicted by $\omega_\pm$;
  \item PlasMol's hybrid curve lies in the same qualitative regime as a
        $\kappa$ estimated from geometry (see Model~2 for a
        geometry-based $\kappa$).
\end{itemize}

% =============================================================================
\section{Script interface}
% =============================================================================

File: \texttt{jobs/model1\_coupled\_oscillators/spectrum\_model1.py}

\begin{itemize}[leftmargin=1.4em]
  \item Edit the \textbf{PARAMETERS} block at the bottom of the script
        ($\omega_p,\gamma_p,f_p,\omega_m,\gamma_m,f_m,\kappa$, grid,
        optional PlasMol CSV overlay).
  \item Run: \texttt{python spectrum\_model1.py}
  \item Outputs: \texttt{spectrum\_model1.csv}, \texttt{spectrum\_model1.png}
\end{itemize}

\noindent
Columns in the CSV:
\texttt{Frequency}, \texttt{Absorption} (normalized),
\texttt{Re\_alpha\_m\_eff}, \texttt{Im\_alpha\_m\_eff}.

% =============================================================================
\section{Limitations}
% =============================================================================

\begin{itemize}[leftmargin=1.4em]
  \item $\kappa$ is phenomenological; it does not know the sphere radius
        or gap.  Use Model~2 for geometry-explicit coupling.
  \item Linear, classical, Markovian damping only---no quantum saturation,
        no non-Markovian continuum of the metal beyond a single Lorentzian
        plasmon mode.
  \item Absolute oscillator strengths $f_p,f_m$ are relative; only ratios
        and $\kappa$ set the hybrid lineshape after peak normalization.
  \item PlasMol's RT-TDDFT Na line and Meep Au dielectric are richer than
        two Lorentzians; agreement is expected at the level of peak count,
        ordering, and rough shifts, not spectroscopic identity.
\end{itemize}

% =============================================================================
\section{References (entry points)}
% =============================================================================

\begin{itemize}[leftmargin=1.4em]
  \item Classical coupled-oscillator / plexciton models for
        plasmon--exciton hybrids (e.g.\ Wu \emph{et al.}, \emph{Nature}
        \textbf{2010}; subsequent plasmon--exciton reviews).
  \item Novotny, ``Strong coupling, energy splitting, and level crossings''
        style treatments of two damped oscillators.
  \item PlasMol Fourier methodology notes
        (\texttt{doc/latex/fourier\_driver\_methodology.tex}): definition
        of $\mu/E_{\mathrm{inc}}$ deconvolution.
\end{itemize}

\end{document}
